\chapter{ON FIELDS}

\textit{July 21st.}

Quizzes will account for 20\% of the final grade. The mid-sem exam will account for 30\% of the final grade,and the final exam will account for 50\%.

\section{An Introduction}

A \eax{field} is a set $F$ with two binary operations, addition and multiplication, satisfying the following axioms:
\begin{enumerate}
    \item $(F,+)$ is an abelian group with (additive) identity element $0$.
    \item $(F\setminus\{0\},\cdot)$ is an abelian group with (multiplicative) identity element $1$.
    \item Multiplication is distributive over addition, i.e., for all $a,b,c\in F$, we have $a\cdot(b+c)=a\cdot b+a\cdot c$.
\end{enumerate}
For a positive integer $n$, define $n \cdot 1 = \underbrace{1 + 1 + \cdots + 1}_{n \text{ times}}$. The \eax{characteristic} of a field $F$ is defined as the smallest positive integer $p$ such that $p \cdot 1 = 0$, if such an $p$ exists; otherwise, the characteristic is defined to be $0$.

\begin{proposition}
    The characteristic of a field is either $0$ or a prime number.
\end{proposition}
\begin{proof}
    Let $F$ be a field with characteristic $\chr F = n > 0$. Let $n = ab$ with $1<a,b<n$. Then $0 = n \cdot 1 = (ab) \cdot 1 = (a \cdot 1)(b \cdot 1)$. Since $F$ is a field, it has no zero divisors, so either $a \cdot 1 = 0$ or $b \cdot 1 = 0$. This contradicts the minimality of $n$, so $n$ must be prime.
\end{proof}

\begin{definition}
    Let $F$ be a field. The \eax{prime subfield} of $F$ is the smallest subfield of $F$ which contains the multiplicative identity $1$.
\end{definition}
This prime subfield can be viewed as both the intersection of all subfields of $F$, and the subfield generated by $1$.
\\

Define, for a positive integer $n$, $(-n) \cdot 1 = - (n \cdot 1)$. Now define $\varphi:\Z \to F$ by $n \mapsto n \cdot 1$ for all $n \in \Z$. It is easy to see that $\varphi$ is a ring homomorphism. It is also routine to check that if $F$ has characteristic $p>0$, then $\ker \varphi = p\Z$, and this induces an injective ring homomorphism $\tilde{\varphi}:\Z/p\Z \to F$. If $F$ has characteristic $0$, then $\ker \varphi = \{0\}$, and $\Z \hookrightarrow F$ via $\varphi$. Since $F$ is a field, it must contain a copy of $\Q$ as well. In the earlier case where $F$ has characteristic $p>0$, then we obtained $\F_{p} = \Z/p\Z \hookrightarrow F$. Hence, for a field of characteristic $0$, the prime subfield is isomorphic to $\Q$, and for a field of characteristic $p>0$, the prime subfield is isomorphic to $\F_{p}$.

\section{Extensions}

Let $F$ be a field. A \eax{field extension} of $F$ is a field $K$ such that $F \subseteq K$. We denote this by $K/F$. Note that $K$ may also be viewed as a vector space over $F$. The \eax{degree} of the extension $K/F$ is defined as the dimension of $K$ as a vector space over $F$, and is denoted by $[K:F]$. If $[K:F]$ is finite, we say that $K/F$ is a \eax{finite extension}. Otherwise, we say that $K/F$ is an \eax{infinite extension}.

The following is the crux of the theory of field extensions; given any polynomial $f(x) \in F[x]$, does there (always) exist an extension of $F$ in which $f(x)$ has a root? Note that it is enough to consider irreducible polynomials, since if $f(x) = g(x)h(x)$, then any root of $g(x)$ or $h(x)$ is a root of $f(x)$. The following theorem answers this question in the affirmative.

\begin{theorem}
    Let $f(x) \in F[x]$ be a non-constant irreducible polynomial. Then there exists a field extension $K/F$ such that $f(x)$ has a root in $K$.
\end{theorem}
\begin{proof}
    Let $K = F[x]/(f(x))$. Since $f(x)$ is irreducible, $(f(x))$ is a maximal ideal of $F[x]$, and hence $K$ is a field. Look at the composition
    \begin{align}
        F \hookrightarrow F[x] \twoheadrightarrow F[x]/(f(x)) = K.
    \end{align}
    This complete composition, $\varphi$, is a field homomorphism, so the kernel must be trivial or the homomorphism is injective. Since $\varphi(1) = 1+(f(x)) \neq 0+(f(x))$, we see that $\ker \varphi = \{0\}$, and hence $F$ is embedded in $K$; $K/F$ is a valid field extension. We claim that $\alpha = x+(f(x)) \in K$ is a root of $f(x)$. Indeed, we have
    \begin{align}
        f(\alpha) = f(x+(f(x))) = f(x) + (f(x)) = 0 + (f(x)) = 0.
    \end{align}
\end{proof}


\begin{example}
    As a concrete example, let $f(x) = x^{2}-2 \in \Q[x]$. Then $f(x)$ is irreducible over $\Q$, and hence there exists a field extension $K/\Q$ such that $f(x)$ has a root in $K$. In fact, we can take $K = \Q[x]/(x^{2}-2)$. Let $\alpha = x+(x^{2}-2) \in K$. Then $\alpha$ is a root of $f(x)$. The reader may also verify that $\{1+(x^{2}-2),\alpha\}$ is a basis of $K$ over $\Q$.
\end{example}
In fact, if $f(x)$ is a polynomial of degree $n$, then the extension $K/F$ constructed above has degree $[K:F] = n$, and a basis of $K$ over $F$ is given by $\{1,\alpha,\alpha^{2},\ldots,\alpha^{n-1}\}$, where $\alpha = x+(f(x)) \in K$. This is exactly the following theorem.

\begin{theorem}
    Let $p(x) \in F[x]$ be an irreducible polynomial of degree $n$, and let $K = F[x]/(p(x))$. Then $[K:F] = n$, and a basis of $K$ over $F$ is given by $\{1,\alpha,\alpha^{2},\ldots,\alpha^{n-1}\}$, where $\alpha = x+(p(x)) \in K$.
\end{theorem}
\begin{proof}
    Note that $K$, as a set, is precisely $\{a_{0}+a_{1}\alpha+\cdots+a_{n-1}\alpha^{n-1} \mid a_{i} \in F\}$, since any polynomial of degree $n$ or higher can be reduced modulo $p(x)$ to a polynomial of degree less than $n$. Hence, the set $\{1,\alpha,\alpha^{2},\ldots,\alpha^{n-1}\}$ spans $K$ over $F$. To see that this set is linearly independent, suppose we have a linear combination
    \begin{align}
        a_{0}+a_{1}\alpha+\cdots+a_{n-1}\alpha^{n-1} = 0
    \end{align}
    Since $p(x)$ is irreducible, the only way for this to happen is if all coefficients $a_{i} = 0$. Thus, the set $\{1,\alpha,\alpha^{2},\ldots,\alpha^{n-1}\}$ is linearly independent.
\end{proof}

\begin{example}
    Consider the polynomial $p(x) = x^{2}+x+1 \in \F_{2}[x]$. This polynomial is irreducible over $\F_{2}$. Let $K = \F_{2}[x]/(x^{2}+x+1) = \{a_{0}+a_{1}\alpha \mid a_{0},a_{1} \in \F_{2}, \, \alpha^{2}+\alpha+1 = 0\}$. This $K$ is an extension of degree $2$ over $\F_{2}$, 
\end{example}

Let $K/F$ be an extension and let $\alpha \in K$. Denote by $F(\alpha)$ the smallest subfield of $K$ containing both $F$ and $\alpha$. We say that $F(\alpha)$ is the subfield of $K$ generated by $\alpha$ over $F$. Of course, we can have multiple generators, and we denote by $F(\alpha_{1},\ldots,\alpha_{n})$ the smallest subfield of $K$ containing $F$ and $\alpha_{1},\ldots,\alpha_{n} \in K$. We say that $F(\alpha_{1},\ldots,\alpha_{n})$ is the subfield of $K$ generated by $\alpha_{1},\ldots,\alpha_{n}$ over $F$. We say $F(\alpha_{1},\alpha_{2},\ldots)$ is \eax{finitely generated} if there exists a finite set of generators $\alpha_{1},\ldots,\alpha_{n}$ such that $F(\alpha_{1},\ldots,\alpha_{n}) = F(\alpha_{1},\alpha_{2},\ldots)$. Otherwise, we say that $F(\alpha_{1},\alpha_{2},\ldots)$ is \eax{infinitely generated}. If the generating set has a single element $\alpha$, we say that $F(\alpha)$ is a \eax{simple extension} of $F$.

\begin{theorem}
    Let $F$ be a field and $p(x)$ an irreducible polynomial over $F$. Suppose $K/F$ is a field extension and $K$ contains a root $\alpha$ of $p(x)$. Then $F(\alpha) \cong F[x]/(p(x))$.
\end{theorem}
\begin{proof}
    Define the map $\varphi : F[x] \to F(\alpha)$ by $f(x) \mapsto f(\alpha)$. This is a ring homomorphism. Clearly, $\varphi|_{F} = \id_{F}$, and $(p(x)) \subseteq \ker \varphi$. Since $p(x)$ is irreducible, $(p(x))$ is a maximal ideal of $F[x]$, and hence $\ker \varphi = (p(x))$. This induces an injective map $\tilde{\varphi} : F[x]/(p(x)) \to F(\alpha)$ given by $f(x)+(p(x)) \mapsto f(\alpha)$. So $\Img \tilde{\varphi}$ is a subfield of $F(\alpha)$ containing both $\alpha$ and $F$. By the very definition of $F(\alpha)$, we must have $\Img \tilde{\varphi} = F(\alpha)$, showing $\varphi$ is surjective. The isomorphism holds.
\end{proof}